\section*{Capítulo \ref{ch:g12:respiration-fermentation} --- A respiração celular e a fermentação}

\begin{solution}{exo:g12:respiration-fermentation:1}
Um \omterm{def:g10:universal-dna:nucleotide}{nucleotídeo} com três fosfatos em fila, a moeda energética da \omterm{def:g10:cells-common-unit:cell}{célula}.
Usá-lo retira o último fosfato, liberando cerca de \qty{30}{kJ} por mol
para impulsionar um processo, e deixa ADP para ser recarregado.
\end{solution}

\begin{solution}{exo:g12:respiration-fermentation:2}
\omterm{prop:g12:respiration-fermentation:glycolysis}{Glicólise} no \omterm{def:g10:cells-common-unit:cell}{citoplasma}: 2 \omterm{def:g12:respiration-fermentation:atp}{ATP} e 2 NAD reduzidos, dois piruvatos. \omterm{prop:g12:respiration-fermentation:mitochondrion}{Ciclo
de Krebs} na matriz mitocondrial: gás carbônico, transportadores
carregados, 2 \omterm{def:g12:respiration-fermentation:atp}{ATP}. Cadeia respiratória na membrana interna:
transportadores descarregados no oxigênio, água, cerca de 26 \omterm{def:g12:respiration-fermentation:atp}{ATP}.
\end{solution}

\begin{solution}{exo:g12:respiration-fermentation:3}
Duas membranas, com a interna dobrada em cristas, envolvendo a matriz.
Matriz: o \omterm{prop:g12:respiration-fermentation:mitochondrion}{ciclo de Krebs}. Membrana interna: a cadeia respiratória e a
síntese de \omterm{def:g12:respiration-fermentation:atp}{ATP}.
\end{solution}

\begin{solution}{exo:g12:respiration-fermentation:4}
Ela descarrega o NAD reduzido da \omterm{prop:g12:respiration-fermentation:glycolysis}{glicólise} no piruvato, para que a
\omterm{prop:g12:respiration-fermentation:glycolysis}{glicólise} consiga continuar sem oxigênio. Rendimento: só os 2 \omterm{def:g12:respiration-fermentation:atp}{ATP} da
\omterm{prop:g12:respiration-fermentation:glycolysis}{glicólise}, mais lactato ou etanol e $\mathrm{CO_2}$.
\end{solution}

\begin{solution}{exo:g12:respiration-fermentation:5}
Cerca de 30 \omterm{def:g12:respiration-fermentation:atp}{ATP} contra 2. A respiração deixa gás carbônico e água; a
\omterm{def:g10:cell-metabolism:fermentation}{fermentação} deixa lactato, ou etanol e gás carbônico, ainda ricos em
energia.
\end{solution}

\begin{solution}{exo:g12:respiration-fermentation:6}
Glicose: efeito nenhum --- as \omterm{def:g10:cells-common-unit:organelle}{mitocôndrias} não conseguem usá-la.
Piruvato: o consumo de oxigênio começa --- o \omterm{def:g11:enzymes-and-phenotype:enzyme}{substrato} da \omterm{def:g10:cells-common-unit:organelle}{mitocôndria}.
ADP: o consumo acelera --- o uso de oxigênio está acoplado à síntese de
\omterm{def:g12:respiration-fermentation:atp}{ATP}. Cianeto: o consumo para --- a cadeia respiratória é o local do uso
de oxigênio.
\end{solution}

\begin{solution}{exo:g12:respiration-fermentation:7}
A \omterm{prop:g12:respiration-fermentation:glycolysis}{glicólise}, que converte glicose em piruvato, ocorre no \omterm{def:g10:cells-common-unit:cell}{citoplasma}; a
\omterm{def:g10:cells-common-unit:organelle}{mitocôndria} não tem as \omterm{def:g11:enzymes-and-phenotype:enzyme}{enzimas} dela e capta piruvato, não glicose.
\end{solution}

\begin{solution}{exo:g12:respiration-fermentation:8}
Cada glicose carrega dois NAD; a \omterm{def:g10:cells-common-unit:cell}{célula} tem apenas um pequeno estoque
de NAD e, uma vez carregado todo ele, a \omterm{prop:g12:respiration-fermentation:glycolysis}{glicólise} não tem transportador
a quem entregar o hidrogênio dela e para. A \omterm{def:g10:cell-metabolism:fermentation}{fermentação} descarrega o
NAD no piruvato e mantém a via rodando.
\end{solution}

\begin{solution}{exo:g12:respiration-fermentation:9}
Respiração: $60/30 = 2$ glicoses por segundo. \omterm{def:g10:cell-metabolism:fermentation}{Fermentação}: $60/2 = 30$
glicoses por segundo, quinze vezes mais.
\end{solution}

\begin{solution}{exo:g12:respiration-fermentation:10}
O cianeto bloqueia a cadeia respiratória, e assim elétron nenhum chega
ao oxigênio e cerca de 26 dos 30 \omterm{def:g12:respiration-fermentation:atp}{ATP} por glicose se perdem na hora; a
\omterm{def:g10:cell-metabolism:fermentation}{fermentação} não consegue cobrir a falta. O cérebro e o coração usam \omterm{def:g12:respiration-fermentation:atp}{ATP}
mais depressa e têm as menores reservas, e assim falham em minutos.
\end{solution}

\begin{solution}{exo:g12:respiration-fermentation:11}
O gás carbônico é liberado na matriz mitocondrial, pelo \omterm{prop:g12:respiration-fermentation:mitochondrion}{ciclo de Krebs}
(e pelo passo que admite o piruvato nele): nenhum na \omterm{prop:g12:respiration-fermentation:glycolysis}{glicólise}. A água
é produzida no fim da cadeia respiratória, quando elétrons e prótons se
juntam ao oxigênio.
\end{solution}

\begin{solution}{exo:g12:respiration-fermentation:12}
O músculo obteve 2 \omterm{def:g12:respiration-fermentation:atp}{ATP} por glicose fermentada; reconstruir essa glicose
a partir do lactato custa ao fígado cerca de 6 \omterm{def:g12:respiration-fermentation:atp}{ATP}, pagos respirando
outro combustível. O corpo gastou 6 para obter 2: a \omterm{def:g10:cell-metabolism:fermentation}{fermentação} toma
emprestada uma energia que precisa ser devolvida com juros.
\end{solution}

\begin{solution}{exo:g12:respiration-fermentation:13}
No ar, a respiração dá 30 \omterm{def:g12:respiration-fermentation:atp}{ATP} por glicose, e assim pouca glicose é
necessária e nenhuma é desviada para etanol; fechada, a \omterm{def:g10:cell-metabolism:fermentation}{fermentação} dá
2, e assim quinze vezes mais glicose é consumida e etanol é produzido.
A massa cresce com gás carbônico, que as duas rotas produzem (a
respiração seis por glicose, a \omterm{def:g10:cell-metabolism:fermentation}{fermentação} dois); na massa densa, de
todo jeito, o oxigênio acaba depressa.
\end{solution}

\begin{solution}{exo:g12:respiration-fermentation:14}
Calor: a energia da cadeia, não mais capturada como \omterm{def:g12:respiration-fermentation:atp}{ATP}, é liberada
diretamente. Um recém-nascido perde calor depressa e não consegue
tremer bem; a gordura marrom é um aquecedor embutido.
\end{solution}

\begin{solution}{exo:g12:respiration-fermentation:15}
Elétrons: a \omterm{def:g10:cell-metabolism:autotrophy}{fotossíntese} os tira da água e os põe no açúcar; a
respiração os tira do açúcar e os dá ao oxigênio. Gases: a \omterm{def:g10:cell-metabolism:autotrophy}{fotossíntese}
consome $\mathrm{CO_2}$ e libera $\mathrm{O_2}$; a respiração, o
contrário. Energia: a \omterm{def:g10:cell-metabolism:autotrophy}{fotossíntese} recebe luz e a armazena no açúcar; a
respiração a tira do açúcar como \omterm{def:g12:respiration-fermentation:atp}{ATP} e calor. \omterm{def:g10:cells-common-unit:organelle}{Organelas}: \omterm{def:g12:photosynthesis:chloroplast}{cloroplasto} e
\omterm{def:g10:cells-common-unit:organelle}{mitocôndria}. Cada uma consome o que a outra produz: a \omterm{def:g10:cell-metabolism:autotrophy}{fotossíntese}
esgotaria o $\mathrm{CO_2}$ e a respiração o $\mathrm{O_2}$ e o alimento.
\end{solution}

\begin{solution}{pb:g12:respiration-fermentation:1}
\textbf{1.} Um terço de \qty{1000}{W} são \qty{300}{W}: \qty{18}{kJ}
por minuto, ou seja, \qty{0.6}{mol} de \omterm{def:g12:respiration-fermentation:atp}{ATP} por minuto.

\textbf{2.} Cerca de \qty{300}{g} por minuto --- contra alguns gramas
guardados: o estoque é reciclado sem parar.

\textbf{3.} \qty{50}{g} são \qty{0.1}{mol}: cada molécula é recarregada
cerca de 6 vezes por minuto.

\textbf{4.} $0.6/30 = \qty{0.02}{mol}$ de glicose por minuto (cerca de
\qty{3.6}{g}).

\textbf{5.} $6 \times 0.02 = \qty{0.12}{mol}$ de $\mathrm{O_2}$, cerca
de \qty{2.9}{L} por minuto --- abaixo do máximo de \qty{4}{L/min} dela
e coerente com os \qty{20}{kJ} por litro de oxigênio do
\cref{ch:g10:exercise-and-energy}: os \qty{1000}{W} conseguem ser
sustentados pela respiração.

\textbf{6.} 2 no \omterm{def:g10:cells-common-unit:cell}{citoplasma}, 28 na \omterm{def:g10:cells-common-unit:organelle}{mitocôndria}: cerca de 93\% do \omterm{def:g12:respiration-fermentation:atp}{ATP}.

\textbf{7.} $6 \times 0.02 = \qty{0.12}{mol}$ de $\mathrm{CO_2}$ por
minuto, todos vindos do \omterm{prop:g12:respiration-fermentation:mitochondrion}{ciclo de Krebs} e da entrada do piruvato na matriz.

\textbf{8.} $30 \times 30 = \qty{900}{kJ}$ dos \qty{2870}{kJ}, cerca de
um terço; o resto é calor, que ela dissipa suando e pelo fluxo de
sangue até a pele dela.

\textbf{9.} As \omterm{def:g11:gene-expression:protein}{proteínas} da cadeia e as \omterm{def:g11:enzymes-and-phenotype:enzyme}{enzimas} que produzem \omterm{def:g12:respiration-fermentation:atp}{ATP} ficam
na membrana; a taxa de síntese de \omterm{def:g12:respiration-fermentation:atp}{ATP} é proporcional à área que as
abriga.

\textbf{10.} Isso eleva a taxa em que a respiração consegue produzir
\omterm{def:g12:respiration-fermentation:atp}{ATP} (e a parcela de gordura que ela consegue queimar), portanto a
potência que ela consegue sustentar sem \omterm{def:g10:cell-metabolism:fermentation}{fermentação}; não muda o
rendimento por glicose nem o oxigênio por glicose.

\textbf{11.} $800 \times 30 = \qty{24}{kJ}$: \qty{0.8}{mol} de \omterm{def:g12:respiration-fermentation:atp}{ATP}.

\textbf{12.} Respiração: $400 \times 30 = \qty{12}{kJ}$, ou seja,
\qty{0.4}{mol}; a \omterm{def:g10:cell-metabolism:fermentation}{fermentação} precisa fornecer os outros \qty{0.4}{mol}.

\textbf{13.} $0.4/2 = \qty{0.2}{mol}$ de glicose, deixando \qty{0.4}{mol}
de lactato.

\textbf{14.} Respiração: $1/30$ de glicose por \omterm{def:g12:respiration-fermentation:atp}{ATP}; \omterm{def:g10:cell-metabolism:fermentation}{fermentação}: $1/2$
--- quinze vezes mais glicose. O músculo aceita isso porque o \omterm{def:g12:respiration-fermentation:atp}{ATP} é
necessário agora e o \omterm{def:g10:exercise-and-energy:glycogen}{glicogênio} está ali; a dívida é paga depois.

\textbf{15.} O lactato precisa ser oxidado ou reconstruído em glicose,
o que custa \omterm{def:g12:respiration-fermentation:atp}{ATP} produzido pela respiração: o consumo de oxigênio fica
acima do repouso até a dívida ser quitada.

\textbf{16.} \qty{500}{W} por uma hora são \qty{1800}{kJ}: $1800/38
\approx \qty{47}{g}$ de gordura.

\textbf{17.} A gordura é mais reduzida --- mais rica em hidrogênio ---,
e assim cada grama precisa de mais oxigênio para ser oxidado por
completo; por litro de oxigênio ela rende um pouco menos energia que a glicose.

\textbf{18.} A \omterm{prop:g12:respiration-fermentation:glycolysis}{glicólise} e a \omterm{def:g10:cell-metabolism:fermentation}{fermentação} dela: a gordura não pode ser
fermentada, e assim um sprint, que roda com \omterm{def:g10:cell-metabolism:fermentation}{fermentação}, não consegue
usar gordura de jeito nenhum --- só glicose e \omterm{def:g10:exercise-and-energy:glycogen}{glicogênio}.

\textbf{19.} O fígado libera glicose do \omterm{def:g10:exercise-and-energy:glycogen}{glicogênio} dele e reconstrói um
pouco a partir do lactato, mantendo a glicose do sangue perto de
\qty{1}{g/L} (\cref{ch:g12:glucose-and-diabetes}); o suprimento do
cérebro fica protegido.

\textbf{20.} Cerca de \qty{0.6}{mol} de \omterm{def:g12:respiration-fermentation:atp}{ATP} por minuto; cerca de
\qty{3}{L} de oxigênio por minuto os pagam; a cadeia respiratória da
membrana interna mitocondrial produz uns nove décimos deles.
\end{solution}
